\section{Introduction}
\label{sec:intro}

A deployed LLM may be asked to remove private, copyrighted, outdated, or
hazardous information long after training. Retraining for every request is
usually impractical, motivating approximate machine unlearning
\cite{cao2015towards,bourtoule2021machine,liu2024rethinking}. Its operational
contract is selective: remove the requested information while preserving
knowledge and capabilities outside the request.

This contract is commonly judged with aggregate forgetting and retention
scores. Even multi-axis evaluations usually summarize a retained population at
the model or checkpoint level. Aggregation makes methods comparable, but it
can hide the distribution of damage. Two checkpoints with similar average
retained utility may differ sharply: one degrades mildly across many behaviors,
whereas the other fails on a small, semantically coherent tail. Fine-grained
audits reveal such localized failures, but typically after unlearning or
inside a neighborhood chosen in advance
\cite{chang2025retain,cheng2026entanglement,ko2025holes}.

This creates a practical allocation problem. A deployer may have thousands of
retained behaviors that could be audited or rehearsed, but enough compute to
protect only a small subset. The useful question is therefore prospective:
for a fixed checkpoint, forget request, and parent unlearner, which retained
behaviors are most at risk once the requested forgetting level is reached, and
where should a limited protection budget be spent?

Recent work shows that collateral damage can be predicted before unlearning
from features of forget--evaluation-set pairs
\cite{su2026preunlearn,su2026corruption}. We study a different resolution and
decision. Rather than predict one outcome for an evaluation set, we rank
individual retained behaviors for the target request from the frozen
pre-unlearning checkpoint. Rather than train a multivariate auditor, we use two
model-derived coordinates and calibrate only a scalar mixture weight on
target-request-disjoint development trajectories.
Finally, we test the rank twice: first as a prospective audit and then as a
fixed-budget repair selector. Prediction correlation alone is not treated as
evidence of decision value.

Our profile separates two questions about each retained behavior. Exact
\emph{Loss Susceptibility} $q_G^\star$ is gradient-based: it is the squared
norm of the behavior loss gradient restricted to a declared parameter block at
the frozen checkpoint. Its forward-only estimate $\widehat q_G$ uses shared
symmetric perturbations and requires no candidate-side backward pass during
target profiling.
\emph{Representation Proximity} $\widehat q_H$ asks how exposed the behavior
is to the current forget request in answer-token hidden-state space. After
rank normalization, a development-calibrated mixture produces one sealed
two-signal score. The two coordinates are complementary descriptors, not
a causal decomposition, and whether their combination improves prediction is
an empirical question.

We organize the paper around three distinct research questions:

\paragraph{RQ1: Prospective prediction.}
Does the sealed two-signal score rank later candidate-level damage and recover
its harmful tail on held-out audit behaviors, improving over each signal alone
and over predeclared simple controls?

\paragraph{RQ2: Validity of Loss Susceptibility estimation.}
On a declared setting-level fidelity support and at a frozen operating point,
does the forward-only estimate recover the ordering and upper-ranked set
induced by exact Loss Susceptibility while avoiding candidate-side
backpropagation during target profiling? The setting-level certificate does
not establish rank fidelity against $q_G^\star$ on every target candidate
support; target profiles separately require outcome-blind numerical
eligibility checks.

\paragraph{RQ3: Decision value.}
Under the same selection quota, repair operator, compute budget, and
feasibility contract, does the preselected two-signal pool improve held-out
mean and tail damage relative to the two single-signal selectors, random
selection, and no repair, without degrading the native retained metric?

The protocol is deliberately sealed. Candidate splits, the operating point for
the Loss Susceptibility estimator, the mixture weight, the repair
quota, and all guards are fixed
before target outcomes are opened. The parent unlearner receives neither the
profile nor the repair pool. Prediction is evaluated only after the parent
reaches a predeclared forgetting gate; protection starts from that same entry
checkpoint and is evaluated on behaviors that never enter repair.

\paragraph{Empirical preview.}
{\setlength{\emergencystretch}{3em}
\PredictionHeadline\ \TailHeadline\ \FidelityHeadline\
\ProtectionHeadline\ \TransferHeadline\par}

Our contributions are:
\begin{itemize}
    \item \textbf{A forward-only target profiler.} We define gradient-based
    Loss Susceptibility, introduce its pool-shared forward-only estimator, and
    combine the estimate with request-specific Representation Proximity.
    \item \textbf{A sealed candidate-level prediction test.} We evaluate one
    pre-unlearning score against later held-out damage, both across the full
    audit ranking and in the harmful tail, with single-signal and simple
    controls on common support.
    \item \textbf{A direct protection test.} We reuse the same sealed score to
    allocate a fixed repair quota and compare selectors under a common
    operator, compute budget, and forgetting/utility contract. This separates
    predictive association from measurable decision value.
\end{itemize}
